% !PRINTED 285
\weberpage{285}

Wenn wir nun, wie es in der Zahlentheorie üblich ist (\gesperrt{Gauss}, Disquisitiones arithmeticae, art. 31), durch das Symbol
\[
\frac{a}{b} \ (\mathrm{mod}\ p)
\]
eine \gesperrt{ganze Zahl} [oder Zahlklasse (mod $p$)] verstehen, die, mit $b$ multipliziert, bei der Teilung durch $p$ den Rest $a$ läßt, so können wir, vorausgesetzt, daß $\mu'$, $\nu'$ nicht durch $p$ teilbar sind, die eine Reihe definierende Kongruenz (2) auch so schreiben:
\weq{3}{\frac{\mu}{\mu'} \equiv \frac{\nu}{\nu'} \ (\mathrm{mod}\ p),}
und wir werden also naturgemäß darauf geführt, durch den Wert des Verhältnisses
\weq{4}{\frac{\mu}{\mu'} \equiv z \ (\mathrm{mod}\ p),}
das jeder der Zahlen $0, 1, \ldots, p-1$ kongruent sein kann, und das für eine ganze Reihe unveränderlich ist, diese Reihe $R$ zu bezeichnen. Es bleibt dabei zunächst die eine Reihe unbezeichnet, in der $\mu'$ und folglich alle $\nu'$ durch $p$ teilbar sind, aber auch diese Reihe ordnet sich der allgemeinen Bezeichnung sehr gut unter, wenn wir, falls $\mu' \equiv 0 \ (\mathrm{mod}\ p)$
\weq{5}{\frac{\mu}{\mu'} \equiv \infty \ (\mathrm{mod}\ p)}
setzen, so daß wir also noch eine $(p+1)$te Reihe $R_\infty$ erhalten. Die Gesamtheit der Reihen, deren Anzahl $p+1$ beträgt, ist hiernach zu bezeichnen durch
\weq{6}{R_\infty, R_0, R_1, \ldots, R_{p-1}.}

Entsprechend werden die zugehörigen Wurzeln einer Transformationsgleichung mit
\weq{7}{v_\infty, v_0, v_1, \ldots, v_{p-1}}
zu bezeichnen sein.

Wenn wir beispielsweise die Invariantengleichung (§ 69) zugrunde legen, so ist [nach den Bestimmungen des § 68, (10) über die Zahlen $a$, $c$, $\partial$]
\weq{8}{v_\infty = j(p\,\omega), \quad v_z = j\left(\frac{z+\omega}{p}\right)}
zu setzen, und ebenso, wenn irgend eine andere Transformationsgleichung gewählt wird.
